% =============================================================================
% Appendix F  The Shared Pipeline: Identities and Formulas
% Cited from sec:pipeline-walk, sec:pipeline-qite, sec:pipeline-readout.
% Plan: ThesisPlanning.md §6 (F); supporting contribution S3 (claims C-03,
% C-04, C-08).
%
% CONTENT  Proofs of prop:stacked-reflection and prop:filter-normalisation;
%          qubitisation (Pi W^k Pi = T_k(H')); controlled-V from a controlled
%          reflection; the filter coefficients; the beta and d formulas; the
%          success probability; the readout identity for p_G < 1.
% SOURCES  Planning.md §3.7 (composition), §3.8, §7.3, §7.6, App. A-7;
%          Journal E023-E024; pihm/src/pihm/circuits/compose.py, gqsp.py,
%          qsp/filters.py, qsp/readout.py.
% =============================================================================

\subsection{Proof of \texorpdfstring{\Cref{prop:stacked-reflection}}{the stacked-residual reflection}}
\label{app:proof-stacked-reflection}
% LAND: V is a Hermitian unitary because it is a unitary conjugation of a
%   reflection; Pi_in V Pi_in = 2 Pi_in U_R^dagger Pi_{a=0} U_R Pi_in - Pi_in, and
%   Pi_in U_R^dagger Pi_{a=0} U_R Pi_in = sum_j R_j^dagger R_j / alpha_R^2 =
%   H / alpha_R^2 by eq:stacked-encoding (the row-select register is left free
%   by Pi_{a=0}). Controlled-V: C(V) = U_R^dagger C(2 Pi - I) U_R, since
%   U_R^dagger U_R = I when the control is off. Numerically: A-7 (1.0e-15).

\subsection{Qubitisation: Chebyshev polynomials of the walk}
\label{app:qubitisation}
% LAND: W acts on each two-dimensional invariant subspace as a rotation by
%   theta = arccos(lambda'), lambda' = 2 lambda / alpha_R^2 - 1; hence
%   Pi_in W^k Pi_in = T_k(H') (eq:walk; lowchuang2019qubitization).
%   CAUTION (lesson L-18): this is ONE matrix power followed by ONE restriction.
%   Iterating "apply W, then project" disagrees outright: the intermediate
%   (ancilla non-zero) branches carry information the final restriction still
%   depends on. Numerically: A-7 (7.7e-15 for k < 12).

\subsection{The controlled walk}
\label{app:controlled-walk}
% LAND: only the two reflections are controlled; U_R and U_R^dagger are not.
%   Why this matters for simulation (lesson L-21): the statevector evolves one
%   controlled walk at a time, memory 16 * 2^q bytes, never a dense operator
%   (4^q). The T1 engine is Qiskit Aer (qiskit2024).

\subsection{The imaginary-time filter}
\label{app:proof-filter-normalisation}
% LAND: the proof of prop:filter-normalisation from the Chebyshev expansion
%   below (the modified-Bessel generating function; abramowitz1964handbook):
%   evaluating at x = -1 gives e^{beta} = I_0 + 2 sum_k I_k(beta), so
%   sum |q_k| = e^{-beta} e^{beta} = 1. Then the truncation degree d and the
%   parameter formulas eq:qite-parameters (the known-gap assumption,
%   lintong2020groundstate), and the success probability p_G -> gamma^2 with
%   amplitude amplification at pi / (4 gamma).
% L-07: scipy's ive returns NaN for large arguments; the coefficients are
%   computed stably (a correctness condition only if the text quotes beta ~ 1e9).

\begin{equation}
    e^{-\beta x} = I_0(\beta) + 2\sum_{k=1}^{\infty} (-1)^k I_k(\beta)\,T_k(x) .
    \label{eq:bessel-expansion}
\end{equation}

\subsection{GQSP phases}
\label{app:gqsp-phases}
% LAND: the circuit R_0 prod_k [C_0 W R_k] (motlagh2024gqsp); the Weiss
%   complement (berntson2025complementary) and the inverse NLFT
%   (laneve2025gqspnlft; CITE-NEEDED arXiv:2505.12615); the realised polynomial
%   against the target on a 4d-point unit-circle grid (<= 1e-12).

\subsection{The readout identity for probabilistic preparation}
\label{app:readout-identity}
% LAND: the all-zero amplitude after U_tau^dagger(x) is <u(x)|v>, with
%   u(x) = tau(x)/sqrt2 the feature map's leading column and u_0 = 1/sqrt(2N);
%   v = sqrt(p_G) psi for the prepared state G and (e_0 + g)/2 for the combined
%   state C. Squaring and subtracting gives eq:readout-identity exactly for any
%   p_G; the paper's Eq. 36 is its p_G = 1 case. Then eq:readout-field.
%   Controlled global phase (lesson L-22): controlled, a circuit's global phase
%   is a relative phase and must become a phase gate on the control; lost, it
%   scales the overlap by cos(lambda).
